\chapter{2015年江苏卷}

\section{填空题}

\begin{problem}
已知集合 \(A = \{1, 2, 3\}\)，\(B = \{2, 4, 5\}\)，则集合 \(A \cup B\) 中元素的个数为 \fillinblank{} \fillinblank{}.
\end{problem}

\begin{problem}
已知一组数据 4, 6, 5, 8, 7, 6，那么这组数据的平均数为 \fillinblank{} \fillinblank{}.
\end{problem}

\begin{problem}
设复数 z 满足 \( z^{2}=3+4\i \) (i 是虚数单位)，则 z 的模为 \fillinblank{} \fillinblank{}.
\end{problem}

\begin{problem}
根据如图所示的伪代码，可知输出的结果为 \fillinblank{}. \[ \begin{array}{l} S \leftarrow 1 \\ I \leftarrow 1 \\ \text { While } I < 8 \\ S \leftarrow S + 2 \\ I \leftarrow I + 3 \\ \text { End While } \\ \text { Print } \end{array} \]
\end{problem}

\begin{problem}
袋中有形状、大小都相同的 4 只球，其中 1 只白球、1 只红球、2 只黄球，从中一次随机摸出 2 只球，则这 2 只球颜色不同的概率为 \fillinblank{} \fillinblank{}.
\end{problem}

\begin{problem}
已知向量 \( a = (2, 1) \) ，\( b = (1, -2) \) ，若 \( ma + nb = (9, -8) \) (m, \( n \in R \) )，则 m - n 的值为 \fillinblank{} \fillinblank{}.
\end{problem}

\begin{problem}
不等式 \(x^{n - 2} < 4\) 的解集为
\end{problem}

\begin{problem}
已知 \(\tan \alpha = -2, \tan (\alpha + \beta) = \frac{1}{7}\)，则 \(\tan \beta\) 的值为 \fillinblank{} \fillinblank{}.
\end{problem}

\begin{problem}
现有橡皮泥制作的底面半径为 5，高为 4 的圆锥和底面半径为 2，高为 8 的圆柱各一个.若将它们重新制作成总体积与高均保持不变，但底面半径相同的新的圆锥和圆柱各一个，则新的底面半径为 \fillinblank{} \fillinblank{}.
\end{problem}

\begin{problem}
在平面直角坐标系 \(xOy\) 中，以点 \((1,0)\) 为圆心且与直线 \(mx - y - 2m - 1 = 0\) (\(m \in \R\)) 相切的所有圆中，半径最大的圆的标准方程为 \fillinblank{} \fillinblank{}.
\end{problem}

\begin{problem}
设数列 \(\{n_n\}\) 满足 \(a_1 = 1\)，且 \(n_{n+1} - a_n = n + 1 (n \in \mathbf{N}^*)\)，则数列 \(\left\{\frac{1}{n_n}\right\}\) 前 10 项的和为 \fillinblank{} \fillinblank{}.
\end{problem}

\begin{problem}
在平面直角坐标系 \(xOy\) 中，\(P\) 为双曲线 \(x^{2} - y^{2} = 1\) 右上方的一个动点，若点 \(P\) 到直线 \(x - y + 1 = 0\) 的距离大于 \(c\) 恒成立，则实数 \(c\) 的最大值为 \fillinblank{} \fillinblank{}.
\end{problem}

\begin{problem}
已知函数 \( f(x) = |\ln x| \)，\( g(x) = \begin{cases} 0, & 0 < x \leqslant 1, \\ |x^2 - 4| - 2, & x > 1, \end{cases} \) 则方程 \( |f(x) + g(x)| = 1 \) 实根的个数为 \fillinblank{} \fillinblank{}.
\end{problem}

\begin{problem}
设向量 \(\mathbf{a}_k = \left(\cos \frac{k\pi}{6}, \sin \frac{k\pi}{6} + \cos \frac{k\pi}{6}\right) (k = 0, 1, 2, \cdots, 12)\)，则 \(\sum_{k=0}^{11} (\mathbf{a}_k \cdot \mathbf{a}_{k+1})\) 的值为 \fillinblank{} \fillinblank{}.
\end{problem}

\section{解答题}

\begin{problem}
在 \(\triangle ABC\) 中，已知 \(AB = 2\)，\(AC = 3\)，\(A = 60^{\circ}\). (1) 求 BC 的长; (2) 求 \( \sin 2C \) 的值.
\end{problem}

\begin{problem}
如图，在直三棱柱 \(ABC - A_{1}B_{1}C_{1}\) 中，已知 \(AC \perp BC\)，\(BC = CC_{1}\)，设 \(AB_{1}\) 的中点为 \(D\)，\(B_{1}C \cap BC_{1} = E\). 求证: (1) \(DE //\) 平面 \(AA_{1}C_{1}C\); (2) \(BC_{1} \perp AB_{1}\).
\end{problem}

\begin{problem}
某山区外围有两条相互垂直的直线型公路，为进一步改善山区的交通现状，计划修建一条连接两条公路和山区边界的直线型公路. 记两条相互垂直的公路为 \(l_{1}, l_{2}\)，山区边界曲线为 \(C\)，计划修建的公路为 \(l\). 如图所示，\(M, N\) 为 \(C\) 的两个端点，测得点 \(M\) 到 \(l_{1}, l_{2}\) 的距离分别为 5 千米和 40 千米，点 \(N\) 到 \(l_{1}, l_{2}\) 的距离分别为 20 千米和 2.5 千米. 以 \(l_{2}, l_{1}\) 所在的直线分别为 \(x, y\) 轴. 建立平面直角坐标系 \(xOy\). 假设曲线 \(C\) 符合函数 \(y = \frac{1}{x^2 + b}\)（其中 \(a, b\) 为常数）模型. (1) 求 \(a, b\) 的值； (2) 设公路 \(l\) 与曲线 \(C\) 相切于 \(P\) 点，\(P\) 的横坐标为 \(t\). \circled{1} 请写出公路 l 长度的函数解析式 \( f(t) \) ，并写出其定义域； \circled{2} 当 t 为何值时，公路 l 的长度最短\(\perp\) 求出最短长度.
\end{problem}

\begin{problem}
如图，在平面直角坐标系 \(xOy\) 中，已知椭圆 \(\frac{x^2}{a^2} + \frac{y^2}{b^2} = 1 (a > b > 0)\) 的离心率为 \(\frac{\sqrt{3}}{2}\)，且右焦点 \(F\) 到左准线 \(l\) 的距离为 3. (1) 求椭圆的标准方程; (2) 过 \(F\) 的直线与椭圆交于 \(A, B\) 两点，线段 \(AB\) 的垂直平分线分别交直线 \(l\) 和 \(AB\) 于点 \(P, C\)，若 \(PC = 2AB\)，求直线 \(AB\) 的方程.
\end{problem}

\begin{problem}
已知函数 \( f(x) = x^3 + ax^2 + b \) (\( a, b \in \R \)). (1) 试讨论 \( f(x) \) 的单调性; (2) 若 \(b = c - a\) (实数 \(c\) 是与 \(a\) 无关的常数). 当函数 \(f(x)\) 有三个不同的零点时，\(a\) 的取值范围恰好是 \((- \infty, -3) \cup \left(1, \frac{3}{2}\right) \cup \left(\frac{3}{2}, +\infty\right)\)，求 \(c\) 的值.
\end{problem}

\begin{problem}
四选二. 【A】如图，在 \(\triangle ABC\) 中， \(AB = AC\) ， \(\triangle ABC\) 的外接圆 \(\odot O\) 的弦 \(AE\) 交\(BC\) 于点 \(D\) .求证： \(\triangle ABD\sim \triangle AEB\). 

【B】已知 \(x, y \in \R\)，向量 \(\alpha = \begin{bmatrix} 1 \\ -1 \\ -2 \end{bmatrix}\) 是矩阵 \(A = \begin{bmatrix} x \\ y \\ 0 \end{bmatrix}\) 的属于特征值 \(-2\) 的一个特征向量，求矩阵 \(A\) 以及它的另一个特征值.
\end{problem}

\begin{problem}
如图，在四棱锥 \(P - ABCD\) 中，已知 \(PA \perp\) 平面 \(ABCD\)，且四边形 \(ABCD\) 为直角梯形，\(\angle ABC = \angle BAD = \frac{\pi}{3}\)，\(PA = AD = 2\)，\(AB = BC = 1\). (1) 求平面 PAB 与平面 PCD 所成二面角的余弦值; (2) 点 Q 是线段 BP 上的动点. 当直线 CQ 与 DP 所成的角最小时，求线段 BQ 的长.
\end{problem}

\begin{problem}
已知集合 \(X = \{1, 2, 3\}\)，\(Y_{n} = \{1, 2, 3, \cdots, n\}\) (\(n \in \mathbf{N}^{*}\))，设 \(S_{n} = \{(a, b) | a\) 整除 \(b\) 或 \(b\) 整除 \(a, a \in X, b \in Y_{n}\}\)，令 \(f(n)\) 表示集合 \(S_{n}\) 所含元素的个数. (1) 写出 \( f(\theta) \) 的值; (2) 当 \(n \geqslant 6\) 时，写出 \(f(n)\) 的表达式，并用数学归纳法证明.
\end{problem}

\begin{problem}
已知集合 \(X = \{1, 2, 3\}\)，\(Y_{n} = \{1, 2, 3, \cdots, n\}\) (\(n \in \mathbf{N}^{*}\))，设 \(S_{n} = \{(a, b) | a\) 整除 \(b\) 或 \(b\) 整除 \(a, a \in X, b \in Y_{n}\}\)，令 \(f(n)\) 表示集合 \(S_{n}\) 所含元素的个数. (1) 写出 \( f(\theta) \) 的值; (2) 当 \(n \geqslant 6\) 时，写出 \(f(n)\) 的表达式，并用数学归纳法证明.
\end{problem}

